\begin{lemma}[The weak-solution identity extends to $W^{1,p}_0$ test functions]
Let $n\ge1$ be an integer and $p>1$ real, and write $p'=\frac p{p-1}$ for the
conjugate exponent.  Let $\Omega\subseteq\R^n$ be open, let
$A\colon\R^n\times\R\times\R^n\to\R^n$, and let $u$, with weak gradient $g$, be a
weak solution of
\[
\operatorname{div}A(x,u,\nabla u)=0\quad\text{in }\Omega
\]
in the sense of \noderef{IsWeakSolutionOn} with right-hand side $B\equiv0$.
Assume
\[
\bigl|A(\cdot,u(\cdot),g(\cdot))\bigr|\in L^{p'}(\Omega).
\]
Then for every $v\in W^{1,p}_0(\Omega)$ with weak gradient $g_v$
(\noderef{MemW1pZeroOn}),
\[
\int_\Omega\bigl(A(x,u(x),g(x)),\,g_v(x)\bigr)\,dx=0 ,
\]
the integrand being the Euclidean inner product
$\sum_{i=1}^n A^i(x,u(x),g(x))\,g_v^i(x)$.
\end{lemma}
\begin{proof}
Throughout write $a(x)=A(x,u(x),g(x))$, so that the hypothesis is
$|a|\in L^{p'}(\Omega)$ and the assertion is
$\int_\Omega(a,g_v)\,dx=0$.

\emph{Step 1: the pairing of $L^{p'}$ with $L^{p}$.}  We record once, for later
use, the following.  Let $F,H\colon\Omega\to\R^n$ be measurable with
$|F|\in L^{p'}(\Omega)$ and $|H|\in L^{p}(\Omega)$.  Then
$x\mapsto(F(x),H(x))$ is integrable on $\Omega$ and
\[
\Bigl|\int_\Omega(F,H)\,dx\Bigr|
\ \le\ \Bigl(\int_\Omega|F|^{p'}dx\Bigr)^{1/p'}
\Bigl(\int_\Omega|H|^{p}dx\Bigr)^{1/p} .
\tag{1}
\]
Indeed, by Cauchy--Schwarz (Preamble) $|(F(x),H(x))|\le|F(x)|\,|H(x)|$
pointwise.  Since $\tfrac1{p'}+\tfrac1p=1$ with $p,p'\in(1,\infty)$, H\"older's
inequality (Preamble) applied to the nonnegative functions $|F|$ and $|H|$ gives
both that $|F|\,|H|$ is integrable on $\Omega$ and that
$\int_\Omega|F||H|\,dx\le\bigl(\int_\Omega|F|^{p'}\bigr)^{1/p'}
\bigl(\int_\Omega|H|^{p}\bigr)^{1/p}$.  Domination (Preamble) then makes
$(F,H)$ integrable, and the norm inequality for the integral together with
monotonicity of the integral (Preamble) gives
$\bigl|\int_\Omega(F,H)\bigr|\le\int_\Omega|(F,H)|\le\int_\Omega|F||H|$, which
is (1).

\emph{Step 2: the approximating sequence.}  By \noderef{MemW1pZeroOn} the
membership $v\in W^{1,p}_0(\Omega)$ with weak gradient $g_v$ provides functions
$\varphi_k\in C^\infty(\R^n)$, $k\in\mathbb N$, each with compact support and
$\operatorname{supp}\varphi_k\subseteq\Omega$, such that
\[
\int_\Omega\Bigl(\sum_{i=1}^n\bigl|g_v^i-D_i\varphi_k\bigr|^2\Bigr)^{p/2}dx
\ \longrightarrow\ 0
\qquad(k\to\infty).
\tag{2}
\]
For a vector $w\in\R^n$ one has $|w|=\bigl(\sum_i|w^i|^2\bigr)^{1/2}$, hence
$\bigl(\sum_i|w^i|^2\bigr)^{p/2}=|w|^{p}$; applying this at
$w=g_v(x)-\nabla\varphi_k(x)$ turns (2) into
\[
\int_\Omega\bigl|g_v-\nabla\varphi_k\bigr|^{p}dx\ \longrightarrow\ 0 .
\tag{2'}
\]
\noderef{MemW1pZeroOn} also gives $v\in W^{1,p}(\Omega)$ with weak gradient
$g_v$, so $|g_v|\in L^{p}(\Omega)$ (\noderef{MemW1pOn}).

\emph{Step 3: each $\varphi_k$ is an admissible test function.}  Fix $k$.  The
function $\varphi_k$ is $C^\infty$ with compact support contained in the open
set $\Omega$, and $1\le n$, $1\le p$; so by
\noderef{SmoothCompactSupportMemW1p}, $\varphi_k\in W^{1,p}(\Omega)$ with weak
gradient $\nabla\varphi_k$ (\noderef{MemW1pOn}).  In particular
$|\nabla\varphi_k|\in L^{p}(\Omega)$.  The pair
$(\varphi_k,\nabla\varphi_k)$ therefore meets the three requirements of the test
class of \noderef{IsWeakSolutionOn} --- membership in $W^{1,p}(\Omega)$ with the
stated weak gradient, compactness of $\operatorname{supp}\varphi_k$, and
$\operatorname{supp}\varphi_k\subseteq\Omega$ --- and the weak-solution
hypothesis gives
\[
\int_\Omega\Bigl[\bigl(a,\nabla\varphi_k\bigr)+B(x,u,g)\,\varphi_k\Bigr]dx=0 .
\]
Here $B\equiv0$, so the second summand vanishes identically and
\[
\int_\Omega\bigl(a,\nabla\varphi_k\bigr)dx=0
\qquad\text{for every }k .
\tag{3}
\]

\emph{Step 4: splitting the integral.}  Fix $k$.  Both $|\nabla\varphi_k|$ and
$|g_v-\nabla\varphi_k|$ lie in $L^{p}(\Omega)$ --- the first by Step 3, the
second because $L^p(\Omega)$ is closed under differences (Minkowski, Preamble)
and $|g_v|\in L^p(\Omega)$ by Step 2.  By Step 1 applied twice, with $F=a$ and
$H=\nabla\varphi_k$ and then with $F=a$ and $H=g_v-\nabla\varphi_k$, both
$x\mapsto(a,\nabla\varphi_k)$ and $x\mapsto(a,g_v-\nabla\varphi_k)$ are
integrable on $\Omega$.  Since
$(a,g_v)=(a,\nabla\varphi_k)+(a,g_v-\nabla\varphi_k)$ pointwise, additivity of
the integral over integrable summands and (3) give
\[
\int_\Omega(a,g_v)\,dx
=\int_\Omega\bigl(a,\nabla\varphi_k\bigr)dx
+\int_\Omega\bigl(a,g_v-\nabla\varphi_k\bigr)dx
=\int_\Omega\bigl(a,g_v-\nabla\varphi_k\bigr)dx .
\tag{4}
\]

\emph{Step 5: passing to the limit.}  The left-hand side of (4) does not depend
on $k$.  Applying (1) with $F=a$ and $H=g_v-\nabla\varphi_k$ to the right-hand
side of (4),
\[
\Bigl|\int_\Omega(a,g_v)\,dx\Bigr|
\ \le\ \Bigl(\int_\Omega|a|^{p'}dx\Bigr)^{1/p'}
\Bigl(\int_\Omega\bigl|g_v-\nabla\varphi_k\bigr|^{p}dx\Bigr)^{1/p}
\qquad\text{for every }k .
\]
The first factor is a finite constant, by the hypothesis $|a|\in L^{p'}(\Omega)$.
By $(2')$ the integral inside the second factor tends to $0$, and $t\mapsto
t^{1/p}$ is continuous at $0$ with value $0$ there because $1/p>0$; so the
right-hand side tends to $0$ as $k\to\infty$.  A constant that is bounded above
by the terms of a sequence converging to $0$ is bounded above by $0$, whence
$\bigl|\int_\Omega(a,g_v)\,dx\bigr|\le0$.  Since an absolute value is
nonnegative, $\bigl|\int_\Omega(a,g_v)\,dx\bigr|=0$ and therefore
$\int_\Omega(a,g_v)\,dx=0$, which is the assertion.
\end{proof}
